% ---------------------------------------------------------------------------
% DRAFT 2 — Introduction, TMLR rebuild
% Spine: strict -> judge-audited -> discriminative -> human (canonical order).
% Repeated facts removed: the Llama-Scout/70B and GPT-OSS 120B/20B comparisons
% now appear once, in Results, not here.
% ---------------------------------------------------------------------------

\section{Introduction}

A benchmark leaderboard reports one number per model, and the field reads that number as a property
of the model. It is not. Accuracy is a joint product of what the model produced and what the scoring
rule was willing to accept, and for open-ended tasks those two things come apart. A model that
computes the right dividend and writes it as \texttt{19,000 crore} where the reference says
\texttt{Rs.\ 19,000 crore} has reasoned correctly and scored zero. Under string-matching evaluation
this is indistinguishable from not knowing the answer.

That string-match and semantic scoring can disagree is not itself a new observation. What this paper
adds is a measurement of how large and structured that disagreement is under controlled
conditions---the same 406 items, the same twelve models, a domain chosen for regulatory text that is
less represented in common Western financial benchmarks---rather than an isolated anecdote: an
eleven-rank reversal for one model, a field-wide
compression to a third of its apparent spread, and an independent finding that roughly half the
items carrying that leaderboard don't discriminate between the systems on it at all.

The usual response is to treat such cases as noise---a small, uniform tax on every model that leaves
rankings intact. We show that this assumption fails, and that it fails hardest for exactly the models
whose output style diverges most from the reference format. Verbose outputs can therefore incur
strict-scoring penalties even when the substantive answer is acceptable, and reasoning-oriented
models that narrate their working are disproportionately exposed to this effect.

Testing this needs a domain less likely to be dominated by memorised recall, so that scoring effects
are not confounded with familiarity. Indian financial regulation supplies
one. SEBI and RBI instruments embed numerical thresholds directly in prose, chain amendments across
decades so that resolving a question requires tracking which circular superseded which, and use
jurisdiction-specific terminology (LODR, PMLA, SFB, AIF, FEMA) that models trained predominantly on
US and European financial corpora have little reason to have learned. Existing financial NLP
benchmarks are built almost entirely from US or European sources (Section~\ref{sec:related}), so
this setting is also unmeasured on its own terms.

We introduce \textbf{IndiaFinBench}\footnote{Dataset, evaluation harness, and all model
predictions: \url{https://anonymous.4open.science/r/IndiaFinBench-322D/}. De-anonymized upon
acceptance.}, 406 expert-annotated items drawn from 34 of 192 collected SEBI and RBI
documents spanning 1992--2026, across four task types: regulatory interpretation, numerical
reasoning, contradiction detection, and temporal reasoning. Annotation quality is validated by a
model-based secondary pass (90.7\% agreement, 150 items) and an independent human
inter-annotator study (180 items, 44.3\% of the benchmark; $\kappa = 0.645$ on the binary
contradiction task). The benchmark is the instrument for the measurement question, and we release
it so the measurement can be repeated elsewhere.

Evaluating twelve contemporary LLMs zero-shot under two scoring regimes produces the paper's central
result. Under a strict four-stage string-matching pipeline, accuracy spans 75.12--89.66\%. Under a
full-coverage cross-model judge audit (Section~\ref{sec:judge})---every REG/NUM/TMP prediction, not
only the ones strict scoring flagged---accuracy spans 94.1--98.0\%. The two regimes do not merely
shift scores; they reorder the leaderboard. DeepSeek-R1-Distill-Llama-70B (hereafter
DeepSeek-R1-Distill) ranks 12th of twelve (last) under strict scoring and 1st under judge-audited
scoring, because 93.9\% of its strict errors are reclassified as correct by the cross-model judge,
primarily reflecting format differences.
Across all twelve models the two rankings are essentially uncorrelated (Spearman $\rho = 0.067$,
$p = 0.84$). A strict leaderboard, in other words, is a poor guide to which model reasons better---
though it remains the right measurement if what you care about is deployment behaviour under a
format-sensitive pipeline. Those are different questions, and accuracy alone does not distinguish
them.

A second measurement problem sits underneath the first. Counting items is not the same as counting
evidence: 213 of the 406 items are answered correctly by at least 11 of the 12 models and therefore
separate almost no pair of systems, leaving 134 items carrying the benchmark's discriminative
signal. We report this for our own benchmark because it is the kind of quantity benchmarks
generally do not report, and because effective size, not nominal size, determines how much a
leaderboard difference can be trusted.

Our contributions, in the order the paper reports them: (1) a two-regime evaluation of twelve LLMs
showing that the scoring rule reorders the leaderboard rather than shifting it uniformly; (2) a
characterisation of who pays the format tax, concentrated in models whose output style diverges from
the reference; (3) an effective-size analysis distinguishing the 134 discriminating items from the
213 that are not; (4) IndiaFinBench itself, with a validated annotation protocol and published
agreement statistics; (5) a human reference point reported with its limits stated; and (6) full
public release of the dataset, harness, per-item judge verdicts, and all model predictions.
